\RequirePackage[thinlines]{easybmat} %-- muss aufgrund eines Bugs vor etex und tikz geladen werden
\documentclass[a4paper, twoside, headsepline, index=totoc,toc=listof, fontsize=10pt, cleardoublepage=empty, headinclude, DIV=12, BCOR=5mm, titlepage]{scrartcl}
\usepackage{scrtime} % KOMA, Uhrzeit ermoeglicht
\usepackage{scrpage2} % wie fancyhdr, nur optimiert auf KOMA-Skript, leicht andere Syntax
\usepackage{etoolbox}
\usepackage{letltxmacro}
\usepackage{ifthen}

%-- Compiler-abhaengige IFs
\usepackage{ifxetex,ifluatex}
\newif\ifxetexorluatex
\ifxetex
  \xetexorluatextrue
\else
  \ifluatex
    \xetexorluatextrue
  \else
    \xetexorluatexfalse
  \fi
\fi

%--Farbdefinitionen
%-- muss vor tikz geladen werden
\usepackage[usenames, table, x11names]{xcolor}
\definecolor{dark_gray}{gray}{0.45}
\definecolor{light_gray}{gray}{0.6}

%--Zum Zeichnen
%-- muss vor polyglossia bzw. babel geladen werden
\usepackage{tikz}
\usepackage{tikz-cd}
\usetikzlibrary{external}
\tikzset{>=latex}
\usetikzlibrary{shapes,arrows,intersections}
\usetikzlibrary{calc,3d}
\usetikzlibrary{decorations.pathreplacing,decorations.markings, decorations.pathmorphing}
\usetikzlibrary{angles}

%-- Konfiguration von tikzexternalize
\tikzexternalize[prefix=tikz/, up to date check=diff]

\ifxetexorluatex
\tikzset{external/system call={xelatex \tikzexternalcheckshellescape %-- verwende LuaLaTeX, wegen dynamischer Speicherallokation
    -halt-on-error -interaction=batchmode -jobname "\image" "\texsource"}}
\else
\tikzset{external/system call={pdflatex \tikzexternalcheckshellescape 
    -halt-on-error -interaction=batchmode -jobname "\image" "\texsource"}}
\fi

%-- tikzexternalize fuer tikzcd deaktivieren, da inkompatibel
\AtBeginEnvironment{tikzcd}{\tikzexternaldisable}
\AtEndEnvironment{tikzcd}{\tikzexternalenable}

%-- um Inkompatibilitaeten von quotes und polyglossia bzw. babel zu vermeiden
\tikzset{
  every picture/.append style={
    execute at begin picture={\shorthandoff{"}},
    execute at end picture={\shorthandon{"}}
  }
}
\usetikzlibrary{quotes}
\usepackage{pgfplots}
\usepgfplotslibrary{colormaps}
\pgfplotsset{compat=1.10}



%-- Mathesymbole etc.
\usepackage{mathtools} % beinhaltet amsmath
\mathtoolsset{showonlyrefs, centercolon}
\newtagform{brackets}[\textbf]{[}{]}
\usetagform{brackets}
\usepackage{wasysym} % zusätzliche Symbole
\usepackage{amssymb} % zusätzliche Symbole
\usepackage{latexsym} % zusätzliche Symbole
\usepackage{stmaryrd} % für Blitz
\usepackage{nicefrac} % schräge Brüche
\usepackage{cancel} % Befehle zum Durchstreichen
\usepackage{mathdots} % Verbesserung von Punkten wie zB \ldots
\usepackage{mathrsfs} % das X wurde benötigt

%-- einzelne Symbole aus dem mathx-package
\DeclareFontFamily{U}{mathx}{\hyphenchar\font45}
\DeclareFontShape{U}{mathx}{m}{n}{<-> mathx10}{}
\DeclareSymbolFont{mathx}{U}{mathx}{m}{n}
\DeclareMathSymbol{\bigplus}{1}{mathx}{"90}

%-- 
\DeclareSymbolFont{bbold}{U}{bbold}{m}{n}
\DeclareSymbolFontAlphabet{\mathbbold}{bbold}
\newcommand{\mathds}[1]{\mathbb{#1}} % Um Kompatibilitaet mit frueheren Benutzung von dsfont herzustellen
\newcommand{\ind}{\mathbbold{1}} % charakteristische-Funktion-Eins

\def\mathul#1#2{\color{#1}\underline{{\color{black}#2}}\color{black}} %farbiges Untersteichen im Mathe-Modus
\renewcommand{\le}{\leqslant}
\renewcommand{\ge}{\geqslant}

%-- Underbrace als Befehl in LaTeX-Syntax (und ohne Spacing probleme mit nachfolgenden Operatoren...)
\newcommand{\Underbrace}[2]{{\underbrace{#1}_{#2}}}

%-- Alles was mit Schrift und XeTeX zu tun hat
\ifxetexorluatex
    % XeLaTeX or LuaTeX
	\ifxetex
		\usepackage{mathspec} %beinhaltet fontspec 
	\else
		\usepackage[no-math]{fontspec}
	\fi
    \usepackage{polyglossia} % moderner babel-ersatz
    \setmainlanguage[spelling=new,babelshorthands=true]{german}
	\shorthandoff{"}
	\setotherlanguage{english}
	% \newcommand\glqq{"}
	% \newcommand\grqq{"}
	\defaultfontfeatures{Mapping=tex-text, WordSpace={1.4}} %
    \setmainfont[Ligatures=Common, BoldFont={* Bold}, ItalicFont={* Light Italic},SmallCapsFont={Linux Libertine O}, SmallCapsFeatures={Letters=SmallCaps}]{Source Sans Pro}
	\setsansfont[Scale=MatchLowercase,Ligatures=Common, BoldFont={* Medium}]{Ubuntu}
	% \newfontface\fancy{texgyreadventor-italic.otf}
	% \setkomafont{title}{\fancy\Huge}
	\ifxetex
		\setallmonofonts[Scale=MatchLowercase, ItalicFont={*}]{Consolas} 
	\else
		\setmonofont[Scale=MatchLowercase, ItalicFont={*}]{Consolas}
	\fi
	\usepackage{xltxtra}
	\usepackage{fontawesome}
	\usepackage{microtype}
\else
    % default: pdfLaTeX
    \usepackage[ngerman]{babel}
    \usepackage[T1]{fontenc}
    \usepackage[utf8]{inputenc}
	\usepackage{textcomp} %verhindert ein paar Fehler bei den Fonts
    \usepackage[babel=true, tracking=true]{microtype}
	\usepackage{lmodern}
	\renewcommand{\familydefault}{\sfdefault} 
\fi


%--Mixed
\usepackage[neverdecrease]{paralist}
\usepackage[german=quotes]{csquotes}
\usepackage{makeidx}
\usepackage{booktabs}
\usepackage{wrapfig}
\usepackage{float}
\usepackage[margin=10pt, font=small, labelfont={sf, bf}, format=plain, indention=1em]{caption}
\captionsetup[wrapfigure]{name=Abb. }
\usepackage{stackrel}
\usepackage{multicol}

\flushbottom


%--Unterstreichung
\usepackage[normalem]{ulem}
\setlength{\ULdepth}{1.8pt}

%--Indexverarbeitung
\newcommand{\bet}[1]{\textbf{#1}}
\newcommand{\Index}[1]{\textbf{#1}\index{#1}}
\makeindex
\setindexpreamble{{\noindent \itshape Die \emph{Seitenzahlen} sind mit Hyperlinks zu den entsprechenden Seiten versehen, also anklickbar} \ifxetexorluatex \faHandUp\fi \par \bigskip}
\renewcommand{\indexpagestyle}{scrheadings}


%--Marginnote & todonotes
\deffootnote[1.5em]{1.5em}{1.5em}{\textsuperscript{\thefootnotemark}\ }
\usepackage[fulladjust]{marginnote}
\renewcommand*{\marginfont}{\color{dark_gray} \itshape \footnotesize}
\usepackage[textsize=small]{todonotes}
\usepackage{ragged2e}
\renewcommand*{\raggedleftmarginnote}{\RaggedLeft}
\renewcommand*{\raggedrightmarginnote}{\RaggedRight}

%--Todonotes mit tikz/externalize kompatibel machen
\LetLtxMacro{\oldtodo}{\todo}
\renewcommand{\todo}[2][]{\tikzexternaldisable\oldtodo[#1]{#2}\tikzexternalenable}
\LetLtxMacro{\oldmissingfigure}{\missingfigure}
\renewcommand{\missingfigure}[2][]{\tikzexternaldisable\oldmissingfigure[{#1}]{#2}\tikzexternalenable}

%--Konfiguration von Hyperref pdfstartview=FitH, 
\usepackage[hidelinks, pdfpagelabels,  bookmarksopen=true, bookmarksnumbered=true, linkcolor=black, urlcolor=SkyBlue2, plainpages=false,pagebackref, citecolor=black, hypertexnames=true, pdfauthor={Jannes Bantje}, pdfborderstyle={/S/U}, linkbordercolor=SkyBlue2, colorlinks=false]{hyperref}

\ifxetexorluatex
\newcommand{\hrefsym}[2]{\href{#1}{\texttt{#2\,\faExternalLink}}}
\newcommand{\hrefsymmail}[2]{\href{#1}{\texttt{\faEnvelope\,#2}}}
\renewcommand{\url}[1]{\hrefsym{#1}{\nolinkurl{#1}}}
\fi
\providecommand{\hrefsym}{\href}
\providecommand{\hrefsymmail}{\hrefsym}


% -- QR-Codes
\usepackage{qrcode} %-- hinter hyperref laden!

%--Römische Zahlen
\newcommand{\RM}[1]{\MakeUppercase{\romannumeral #1{}}}
\renewcommand{\thesection}{\Roman{section}}

%%--Abkürzungen etc.
\newcommand{\light}{\text{\Large $\lightning$}}

%-- Definitionen von weiteren Mathe-Befehlen
\DeclareMathOperator{\re}{Re} %Realteil
\DeclareMathOperator{\im}{Im} %Imaginaerteil
\DeclareMathOperator{\id}{id} %identische Abbildung
\DeclareMathOperator{\Hom}{Hom} 
\DeclareMathOperator{\Mat}{Mat}
\DeclareMathOperator{\Eig}{Eig}



\newcommand{\mathd}{\ensuremath{\mathrm{d}\mkern-1.0mu}}



%--Beweisende
\newcommand{\bewende}{\ifmmode \tag*{$\square$} \else \hfill $\square$ \fi}


%--Skalarprodukt
\DeclarePairedDelimiterX\sprod[2]{\langle}{\rangle}{#1\,\delimsize\vert\,#2}
\DeclarePairedDelimiterX\skal[2]{\langle}{\rangle}{#1\,,\,#2}

%--Betrag, Gaußklammer
\DeclarePairedDelimiter{\abs}{\lvert}{\rvert}
\DeclarePairedDelimiter{\floor}{\lfloor}{\rfloor}
\DeclarePairedDelimiter{\ceil}{\lceil}{\rceil}

%--Norm
\DeclarePairedDelimiter\doppelstrich{\Vert}{\Vert}
\newcommand{\norm}[2][\relax]{
\ifx#1\relax \ensuremath{\doppelstrich*{#2}}
\else \ensuremath{\doppelstrich*{#2}_{#1}}
\fi}


%--Umklammern
\DeclarePairedDelimiter\enbrace{(}{)}
\DeclarePairedDelimiter\benbrace{[}{]}
\DeclarePairedDelimiter\lenbrace{<}{>}
\newcommand{\ssbrace}[1]{{\scriptscriptstyle\enbrace{#1}}}

%--Mengen
\DeclarePairedDelimiterX\mengenA[1]{\lbrace}{\rbrace}{#1}
\DeclarePairedDelimiterX\mengenB[2]{\lbrace}{\rbrace}{#1\, \delimsize\vert \, #2}

\makeatletter
\newcommand{\set}[2][\relax]{
\ifx#1\relax \ensuremath{
\mengenA*{#2}}
\else \ensuremath{%
  \mengenB*{#1}{#2}}
\fi}
\makeatother

\DeclareRobustCommand{\minwidthbox}[2]{%
  \ifmmode
    \expandafter\mathmakebox
  \else
    \expandafter\makebox
  \fi
  [\ifdim#2<\width\width\else#2\fi]{#1}%
}


%--offen und abgeschlossen
\newcommand{\off}{\! \stackrel[\text{offen}]{}{\subset} \!}
\newcommand{\abg}{\! \stackrel[\text{abg.}]{}{\subset} \!}

%--Differential
\newcommand{\diff}[2]{\ensuremath{\frac{{\partial #1}}{{\partial #2}} }}
\newcommand{\diffd}[2]{\ensuremath{\frac{\mathd #1}{\mathd #2} }}

%--Diagonale Linien für Matrizen
\newcommand\tikzmark[1]{%
  \tikzexternaldisable\tikz[overlay,remember picture,baseline] \node [anchor=base] (#1) {};\tikzexternalenable}

\newcommand\MyLine[3][]{%
  \tikzexternaldisable\begin{tikzpicture}[overlay,remember picture]
    \draw[#1] (#2.south) -- (#3.north);
  \end{tikzpicture}\tikzexternalenable}
  
%--Jordankasten
\newcommand{\Jord}{ \tikzexternaldisable \begin{BMAT}{|ccc|}{|ccc|} %
				0 & & \\ %
				1\tikzmark{z} &  \diagdown & \\ %
				& \tikzmark{u}1 \MyLine[thick]{z}{u} &  0 %
			\end{BMAT} \tikzexternalenable}
			
%--Jordankasten mit Lambda
\newcommand{\JordLa}[1][\empty]{%
	\tikzexternaldisable\ifthenelse{\equal{#1}{\empty}}
		{\begin{BMAT}{|ccc|}{|ccc|} %
				\lambda & & \\ %
				1\tikzmark{z} &  \diagdown & \\ %
				& \tikzmark{u}1 \MyLine[thick]{z}{u} &  \lambda  %
			\end{BMAT}}
		{\begin{BMAT}{|ccc|}{|ccc|} %
				\lambda_{#1}  & & \\ %
				1\tikzmark{z} &  \diagdown & \\ %
				& \tikzmark{u}1 \MyLine[thick]{z}{u} &  \lambda_{#1}  %
			\end{BMAT}}\tikzexternalenable}

%--Inhaltsverzeichnis
\usepackage[tocindentauto]{tocstyle}
\usetocstyle{KOMAlike}	
\shorthandon{"}
		
%--Punkte (als letztes laden)
\usepackage{ellipsis}

\newcommand{\Semester}{SoSe 2015}
\newcommand{\fach}{Analysis 2}
\newcommand{\prof}{Prof.\ Ebert}

\providecommand{\verfasser}{Tim Keil}

%--Konfiguration von scrheadings
\setheadsepline{1pt}[\color{light_gray}]
\pagestyle{scrheadings}
\clearscrheadfoot

\providecommand{\shortFach}{\fach}
\lehead{ \includegraphics[height=0.6 cm,keepaspectratio]{../!config/Bilder/Logo_WWU_Muenster_light_gray.pdf}}
\rehead{\rule{0cm}{0.6cm}\footnotesize \sffamily \color{light_gray} \verfasser{} -- Mitschrift \shortFach}
\lohead{\rule{0cm}{0.6cm} \footnotesize \sffamily \color{light_gray} Stand: \today \; \thistime[:]}
\rohead{\includegraphics[height=0.6 cm,keepaspectratio]{../!config/Bilder/fb10logo_gray.pdf}}


\ofoot[{ \color{dark_gray} \LARGE \sffamily \thepage}]{{ \color{dark_gray} \LARGE \sffamily \thepage}} %hier wir auch der plain Stil bearbeitet!
\automark{section}
\ifoot{ \color{dark_gray} \small \leftmark}

%--Metadaten
\providecommand{\mail}{keil.menden@web.de}
\author{\verfasser}
\titlehead{\includegraphics[height=1.5cm, keepaspectratio]{../!config/Bilder/Logo_WWU_Muenster.pdf}%
\hfill \includegraphics[height=1.3cm, keepaspectratio]{../!config/Bilder/fb10logo.pdf}}
\title{Skript \fach}
\subtitle{Mitschrift der Vorlesung  \enquote{\fach} von \prof}



\begin{document}
\maketitle
\pagenumbering{Alph}
\section*{Vorwort --- Mitarbeit am Skript}
Dieses Dokument ist eine Mitschrift aus der Vorlesung \enquote{\fach, \Semester}, gelesen von \prof. Der Inhalt entspricht weitestgehend dem Tafelanschrieb. Für die
Korrektheit des Inhalts übernehme ich keinerlei Garantie! Für Bemerkungen und Korrekturen -- und seien es nur Rechtschreibfehler -- bin ich sehr dankbar. 
Korrekturen bitte durch persönliches Ansprechen oder per Mail an keil.menden@web.de.
%\begin{itemize}
%	\item Persönliches Ansprechen in der Uni oder Mails an keil.menden@web.de (gerne auch mit annotieren PDFs) \url{https://github.com/JaMeZ-B/latex-wwu}.
%	\item \emph{Direktes} Mitarbeiten am Skript: Den Quellcode poste ich auf GitHub (siehe oben), also stehen vielfältige Möglichkeiten der Zusammenarbeit zur Verfügung:
%	Zum Beispiel durch Kommentare am Code über die Website und die Kombination Fork + Pull Request. Wer sich verdient macht oder ein Skript zu einer Vorlesung, die 
%	ich nicht besuche, beisteuern will, dem gewähre ich gerne auch Schreibzugriff.
	
%	Beachten sollte man dabei, dass dazu ein Account bei \url{github.com} notwendig ist, der allerdings ohne Angabe von persönlichen Daten angelegt werden kann. 
%	Wer bei GitHub (bzw. dem zugrunde liegenden Open-Source-Programm \enquote{\texttt{git}}) -- verständlicherweise -- Hilfe beim Einstieg braucht, dem helfe ich gerne 
%	weiter. Es gibt aber auch zahlreiche empfehlenswerte Tutorials im Internet.\footnote{zB. \url{https://try.github.io/levels/1/challenges/1}, ist auf Englisch, aber dafür 
%	interaktives LearningByDoing}
%	\item \emph{Indirektes} Mitarbeiten: \TeX-Dateien per Mail verschicken. 
	
%	Dies ist nur dann sinnvoll, wenn man einen ganzen Abschnitt ändern möchte (zB. einen alternativen Beweis geben), da ich die Änderungen dann per Hand einbauen muss! Ich freue mich aber auch über solche Beiträge!
%\end{itemize}

\newpage

\tableofcontents
\cleardoubleoddemptypage
\pagenumbering{arabic}
\setcounter{page}{1}


\section{Normierte und metrische Räume} % (fold)
\label{sec:Normierte und metrische Räume}

\begin{centering}
	$\mathbb{R} \hspace{1cm} \rightarrow{} \hspace{1cm}  \mathbb{R}^n \hspace{1cm}  \rightarrow{} \hspace{1cm}  $Normierte Vektorräume$ \hspace{1cm} \rightarrow \hspace{1cm}  $Metrische Räume
\end{centering}

\subsection{Der $\mathbb{R}^n$} % (fold)
\label{sub:der_mathbb_r_n}

\[
	\mathbb{R}^n = \left\{ \left. \begin{pmatrix} x_1 \\ \vdots \\ x_n \\ \end{pmatrix}  \right| x_i \in \mathbb{R} \right\}
\]

Der $\mathbb{R}^n$ ist ein $\mathbb{R}$-Vektorraum und $\dim_{\mathbb{R}}(\mathbb{R}^n)=n$.

Sei
\[
	e_i = \begin{pmatrix} 0 \\ \vdots \\ 0 \\ 1 \\ 0 \\ \vdots \\ 0 \end{pmatrix} \leftarrow \text{$i$-ter Eintrag}
\]

\((e_1,\dots,e_n)\) ist eine Basis von $\mathbb{R}^n$ und heißt {\bfseries Standardbasis}. 

Ist \(x = \begin{pmatrix} x_1 \\ \vdots \\ x_n \end{pmatrix}\), so gilt \(x=\sum\limits^{n}_{i=1}x_ie_i\)

\subsubsection[Das Skalarprodukt]{Definition} % (fold)
\label{ssub:definition}

Das {\bfseries Skalarprodukt} auf $\mathbb{R}^n$ ist die Abbildung \[ \langle\, .\,,\,.\,\rangle : \mathbb{R}^n \times \, \mathbb{R}^n \to \mathbb{R}\] mit 

\[
 x = \left< \begin{pmatrix} x_1 \\ \vdots \\ x_n \end{pmatrix} , \begin{pmatrix} y_1 \\ \vdots \\ y_n \end{pmatrix} \right> := \sum_{i=1}^{n}x_iy_i \in \mathbb{R} 
 \]
 
es gilt:
\begin{description}
	\item[(S1)] (Bilinearität) $\forall \, x_0,x_1,y_0,y_1, \in \mathbb{R}^n \hspace{4pt} \forall \, a,b \in \mathbb{R}$
	 \begin{eqnarray}
	 \langle x_0+x_1\,,\, y_0+y_1\rangle & =& \langle x_0\,,\, y_0\rangle + \langle x_1\,,\,y_1\rangle + \langle x_1\,,\, y_0\rangle + \langle x_1\,,\, y_1\rangle \\
	 \langle ax_0\,,\, by_0\rangle &=& ab \langle x_0\,,\, y_0\rangle 
	 \end{eqnarray}
	 \item[(S2)](Symmetrie) $\forall \, x,y \in \mathbb{R}^n$ gilt $\skal{x}{y} = \skal{y}{x}$
	
	\item[(S3)](positive Definitheit) $\forall \, x \in \mathbb{R}$ gilt $\skal{x}{x} \geq 0$ und $\skal{x}{x}=0$ nur dann, wenn $x=0$
\end{description}
Desweiteren gilt $\skal{e_i}{e_j}=\delta_y := \begin{cases} 1 & i=j \\ 0 &i\neq j \end{cases}$ \hspace{1cm} (Kronecker-Symbol)

Ist \[
	x=\sum_{i=1}^{n}x_ie_i \in \mathbb{R}^n \qquad (x_i \in \mathbb{R})
	\]
so ist \begin{gather*}
	\skal{x}{e_j} = \skal{\sum_{i=1}^{n}x_ie_i}{e_j} \stackrel{(S1)}{=} \sum_{i=1}^{n}x_i \skal{e_i}{e_j} = \sum^{n}_{i=1}x_i \delta_{ij} = x_j \\
	\Rightarrow x = \sum^{n}_{j=1} \skal{x}{e_j}e_j
\end{gather*}

\subsubsection[Die $l^2$-Norm]{Definition} % (fold)
\label{ssub:definition}
Sei $x \in \mathbb{R}^n$. Die $l^2$-Norm von $x$ ist die Zahl
\[
	\norm{x}_2:= \sqrt{\skal{x}{x}} 
\]
Schreibweise manchmal auch: $\norm{x}_2 = \norm{x}$
% subsubsection definition (end)

\subsection{Euklidsche und normierte Vektorräume} % (fold)
\label{sub:euklidsche_und_normierte_vektorraume}

\subsubsection[Euklidscher Vektorraum]{Definition} % (fold)
\label{ssub:definition}
Sei $ V$  ein $\mathbb{R}$-Vektorraum. Ein {\bfseries Skalarprodukt} auf $ V$  ist eine Abbildung $\skal{.}{.}:V \times V \to \mathbb{R}$, welche die Axiome S1,S2,S3 erfüllt. \\
Ein {\bfseries euklidscher Vektorraum} ist ein Paar $(V,\skal{}{})$. $ V$  ist ein $\mathbb{R}$-Vektorraum und $\skal{}{}$ ist Skalarprodukt auf $ V$. \\
Ist $(V,\skal{}{})$ ein eukdlischer Vektorraum, so ist die durch $\skal{}{}$ induzierte Norm definiert als \[
	\norm{v}=\sqrt{\skal{v}{v}} \hspace{1cm}  v \in V 
	\]

\subsubsection[Beispiel euklidscher Vektorraum]{Beispiele} % (fold)
\label{ssub:beispiele}

\begin{itemize}
	\item $\mathbb{R}^n$ mit Standardskalarprodukt
	\item Die Abbildung \begin{align*}
		 \mathbb{C} \times \mathbb{C} &\to \mathbb{R} \\
		 (z,w) &\mapsto Re(\bar z \cdot w)
		 \end{align*}
		 ist ein Skalarprodukt. Die induzierte Norm ist der Absolutbetrag.      
\end{itemize}
% subsubsection beispiele (end)

\subsubsection[Norm]{Definition} % (fold)
\label{ssub:definition}
Sei $ V$ ein $\mathbb{R} $-Vektorraum. Eine {\bfseries Norm} auf $ V$ ist eine Abbildung $\norm{.}:V \to \mathbb{R} , v \mapsto \norm{v}$ mit
\begin{description}
	\item[(N1)]$\norm{v} \geq 0 \hspace{1cm} \forall \, v \in V$ 
	\item[(N2)] $\norm{v}=0 \Leftrightarrow v=0$
	\item[(N3)] $\norm{av}=|a|\cdot \norm{v} \hspace{1cm}  \forall\,v \in V, a \in \mathbb{R}$
	\item[(N4)] $\norm{v+w} \leq \norm{v}+\norm{w} \hspace{1cm} \forall\, v,w \in V$ \hspace{1cm} (Dreiecksungleichung) 
\end{description} 
Ein normierter Vektorraum ist ein Paar $(V,\norm{.})$ bestehend aus einem $\mathbb{R}$-VR $ V$ und einer Norm auf $ V$.

\subsubsection[Cauchy-Schwarz-Ungleichung]{Satz} % (fold)
\label{ssub:satz}
Sei $(V,\skal{}{})$ ein euklidscher Vektorraum. Für $v \in V$ setzte $\norm{v}=\sqrt{\skal{v}{v}}$ \\ Dann gilt:
\begin{enumerate}
	\item $\norm{.}$ ist eine Norm auf $ V$ 
	\item $\forall\, v,w \in V$ gilt die {\bfseries Cauchy-Schwarz-Ungleichung}:
	\[
		|\skal{v}{w}| \leq \norm{v} \cdot \norm{w}
	\]  
\end{enumerate}
\minisec{Beweis:}
\begin{enumerate}
	\item \begin{enumerate}[N1:]
		\item $\norm{v}=\sqrt{\underset{\geq 0}{\underbrace{\skal{v}{v}}}} \geq  0$
		\item $\norm{v}=\sqrt{\skal{v}{v}}=0 \Rightarrow \skal{v}{v}=0 \stackrel{\text{S3}}{\Rightarrow} v=0$ 
		\item $a \in \mathbb{R} , v \in V$ \[
			\norm{av}=\sqrt{\skal{av}{av}} \stackrel{\text{S1}}{=} \sqrt{a^2 \skal{v}{v}} = \sqrt{a^2} \sqrt{\skal{v}{v}} = \abs{a} \norm{v}
		\] 
	\end{enumerate}
	\item CS-Ungleichung ist trivial, wenn $w=0$ \\
	Ann.: $w \neq 0$: Sei $f:\mathbb{R} \to \mathbb{R} $ mit
	\begin{align*}
		f(t) &=\skal{v+tw}{v+tw} \\
		& = \skal{v}{v} +2t \skal{v}{w} + t^2 \skal{w}{w} \\
		& = \norm{v}^2+ 2t \underbrace{\skal{v}{w}}_{\in \mathbb{R}} + t^2 \norm{w}^2 
	\end{align*} 
	aus S3 folgt $f(t)\geq 0$. f ist Polynom
	\begin{align*}
		&f'(t)=0 \Leftrightarrow t_0 = - \frac{\skal{v}{w}}{\norm{w}^2} \\
		&f''(t)=2 \norm{w}^2 > 0
	\end{align*}
 $\Rightarrow$ $f$ nimmt bei $t_0$ globales Minimum an
 \begin{gather*}
 	0 \leq f(t_0)=\norm{v}^2 - \frac{2 \skal{v}{w}^2}{\norm{w}^2} + \frac{\skal{v}{w}^2}{\norm{w}^4} \norm{w}^2 = \norm{v}^2 - \frac{\skal{v}{w}^2}{\norm{w}^2} \\
	\Rightarrow \frac{\skal{v}{w}^2}{\norm{w}^2} \leq \norm{v}^2 \\ \Rightarrow \skal{v}{w}^2 \leq \norm{v}^2 \norm{w}^2 \\
	\Rightarrow \abs{\skal{v}{w}} = \sqrt{\skal{v}{w}^2} \leq \sqrt{\norm{v}^2 \norm{w}^2} = \norm{v} \cdot \norm{w}
 \end{gather*}
 \bewende
\end{enumerate}
Schließlich könnnen wir N4 zeigen:
\begin{enumerate}
	\item[N4:]\begin{align*}
		\norm{v+w}^2 &= \skal{v+w}{v+w} \\ 
		&\stackrel{\text{S1}}{=} \norm{v}^2+ 2 \skal{v}{w}+ \norm{w}^2 \\
		&\leq \norm{v} + 2 \abs{\skal{v}{w}}+ \norm{w}^2 \\
		&\stackrel{\text{CS}}{\leq} \norm{v}^2+2 \norm{v}\norm{w}+ \norm{w}^2 \\
		&=(\norm{v}+\norm{w})^2
	\end{align*} 
Wurzel ziehen zeigt N4
\bewende
\end{enumerate}
% subsubsection satz (end)

\subsubsection[Besipiele für Normen auf $\mathbb{R}$]{Weitere Beispiele für Normen auf $\mathbb{R}$} % (fold)
\label{ssub:weitere_beispiele_fur_normen_auf_mathbb_r}

\begin{gather*}
	 x = \sum^{n}_{i=1}x_ie_i \hspace{1cm}  x_i \in  \mathbb{R}
	 \end{gather*}
\begin{itemize}
	\item \text{$l^1$-Norm} $\hspace{1cm}  \norm{x}_1 = \sum^{n}_{i=1}\abs{x_i}$ 
	\item $l^{\infty}$-Norm $\hspace{1cm} \norm{x}_{\infty}=\max{\{\abs{x_i}\,|\,i=1,\dots,n}\}$
	\item Ist $p \in [1,\infty)$, so ist die $l^p$-Norm
	\[
		\norm{x}_p = \left( \sum^{n}_{i=1}\abs{x_i}^p \right)^{\frac{1}{p}} \hspace{1cm} \text{(Beweis Schwer)}
	\] 
\end{itemize}
% subsubsection weitere_beispiele_fur_normen_auf_mathbb_r (end)

\subsubsection[Äquivalenz von Normen]{Definition} % (fold)
\label{ssub:definition}
$ V$ $\mathbb{R}$-VR , $\norm{.}_0 , \norm{.}_1$ Normen auf $ V$. Die Normen $\norm{.}_0 , \norm{.}_1$ heißen {\bfseries äquivalent}, falls $C,c>0$ existieren, so dass für alle $v \in  V$ gilt:
\[
	c \norm{v}_0 \leq \norm{v}_1 \leq C \norm{v}_0 
\] 
\minisec{Bemerkung:} Das ist eine Äquivalenzrelation
\begin{itemize}
	\item {\bfseries Reflexiv}: klar
	\item {\bfseries Symmetrie}: \begin{gather*}
		c \norm{v}_0 \leq \norm{v}_1 \leq C \norm{v}_0 \\
		\Rightarrow \frac{1}{c} \norm{v}_1 \leq \norm{v}_0 \leq \frac{1}{C} \norm{v}_1
	\end{gather*}
	\item {\bfseries Transitiv}: \begin{gather*}
		c \norm{v}_0 \leq \norm{v}_1 \leq C \norm{v}_0 \hspace{0.5cm} \text{und}\hspace{0.5cm} d \norm{v}_1 \leq \norm{v}_2 \leq D \norm{v}_1  \hspace{1cm} (C,c,D,d > 0)	\end{gather*}
		\[
		\Rightarrow cd \norm{v}_0 \leq CD \norm{v}_0
		\]
\end{itemize}

\subsubsection[Beispiel Äquivalente Normen]{Beispiel} % (fold)
\label{ssub:beispiel}
Die $l^1, l^2 \, \text{und} \, l^{\infty}$-Norm auf $\mathbb{R}^n $ sind äquivalent und zwar gilt:
\[
	\norm{v}_{\infty} \leq \norm{v}_2 \leq \norm{v}_1 \leq n \norm{v}_{\infty} \hspace{1cm} \text{(Beweis Übung)}
\]
% subsubsection beispiel (end)

\subsection{Metrische Räume} % (fold)
\label{sub:metrische_raume}

\subsubsection[Abstand]{Definition} % (fold)
\label{ssub:definition}
$(V,\norm{.})$ normierter Vektorraum, $v,w \in  V$ \\
Der {\bfseries Abstand} von $v$ und $w$ ist \[
d(v,w):= \norm{v-w}	
\] 
% subsubsection definition (end)
 \subsubsection[Metrik]{Definition} % (fold)
 \label{ssub:definition}
Sei $X$ eine Menge \\
Eine {\bfseries Metrik} (oder Abstandsfunktion) ist eine Abbildung $d: X \times X \to \mathbb{R} , (x,y) \mapsto d(x,y)$, sodass gilt:
\begin{description}
	\item[(M1)] $\forall\, x,y \in X: \hspace{1cm}  d(x,y) \geq  0$
	\item[(M2)] $\forall\, x,y \in X: \hspace{1cm}  d(x,x)=0   \hspace{1cm} \text{und} \hspace{1cm} 
		 		d(x,y)=0 \Rightarrow x=y$
	\item[(M3)] $\forall\, x,y \in X: \hspace{1cm} d(x,y)=d(y,x)$
	\item[(M4)] $\forall\, x,y,z \in X: \hspace{1cm}  d(x,z) \leq d(x,y)+d(y,z) \hspace{1cm}$ ({\bfseries Dreiecksungleichung}) 
\end{description}
Ein {\bfseries Metrischer Raum} ist ein Paar $(X,d)$, $X$ Menge, $d$ Metrik auf $X$
\minisec{Beispiele:}
\begin{itemize}
	\item Ist $(V,\norm{.})$ normierter Vektorraum, so ist durch $d(v,w):= \norm{v-w}$ eine Metrik auf $ V$ gegeben.

\begin{align*}
	v,w,u \in V \hspace{1cm} d(v,u) &\stackrel{\hphantom{\text{$\triangle$-Ungl.}}}{=}  \norm{v-u} = \norm{v-w+w-u} \\
	&\stackrel{\text{$\triangle$-Ungl.}}{\leq} \norm{v-w} + \norm{w-u} \\ 
	&\stackrel{\hphantom{\text{$\triangle$-Ungl.}}}{=} d(v,w)+d(w,u)
\end{align*}
 
\minisec{Unterbeispiel:}
$V=\mathbb{R}$. Dann ist $x \mapsto \abs{x} $ eine Norm auf $\mathbb{R}$ und die induzierte Metrik auf $\mathbb{R}$ ist $d(x,y)=\norm{x-y}$
 % subsubsection definition (end)
 \item $(X,d_X)$ metrischer Raum $Y\subseteq X$ Teilmenge. Dann ist $(Y,d_y)$ ein metrischer Raum, wobei für $y_0,y_1 \in Y \text{gilt, dass} \, d_y(y_0,y_1):=d_x(y_0,y_1)$ \\
 \underline{Vereinbarung:} Ist $X \subseteq \mathbb{R}^n$ eine Teilmenge, so ist $X$ immer mit der durch das Standardskalarprodukt induzierte Metrik versehen. (es sei denn, etwas anderes ist gesagt) 
\end{itemize}

\minisec{Konvergenz}
Erinnerung (Ana I): $(x_n)_n$ Folge in $\mathbb{R}$, $x \in  \mathbb{R}$. Dann 
\[
\lim_{n \to \infty}x_n=x \, :\Leftrightarrow \, \forall\, \varepsilon > 0 \,\exists \,n_0 \, \forall\, n \geq n_0: \, \abs{x-x_1}< \varepsilon
 \] 
% subsection metrische_raume (end)
\subsubsection[Konvergenz im metrischen Raum]{Definition} % (fold)
\label{ssub:definition}
$(X,d)$ metrischer Raum, $(x_n)_n$ Folge in $X$ (d.h eine Abbildung $x: \mathbb{N} \to X$) $y \in X$. Dann Konvergiert die Folge $(x_n)_n$ gegen $y :\Leftrightarrow 
 \forall\, \varepsilon >0 \, \exists \, n_0: \forall \ n \geq n_0 \, d(x_n,y)< \varepsilon $ \\
\underline{Äquivalent}: Die Folge (in $\mathbb{R}$) $d(x_n,y)$ ist eine Nullfolge.
% subsubsection definition (end)

\subsubsection[Beispiel Metrik]{Beispiel} % (fold)
\label{ssub:3.4}
$X=\mathbb{R}^n$ mit Metrik, die durch $l^1$-$ l^2$-$ \text{ oder } l^{\infty}$- Norm induziert wird. Folge $k \mapsto x(k) \in \mathbb{R}^n$.
\[
	x(k)= \sum_{i=1}^{n}x_i(k)e_i \qquad e_i \text{ Standard-Basisvektor} \qquad x_i(k) \in  \mathbb{R} \qquad x=\sum^{n}_{i=1}x_ie_i
\]
Dann sind äquivalent
\begin{enumerate}
	\item für alle $i=1,\dots,n$ gilt $\lim\limits_{k \to \infty}x_i(k)=x_i$
	\item $\lim\limits_{k \to \infty}\norm{x(k)-x}_1=0$
	\item $\lim\limits_{k \to \infty}\norm{x(k)-x}_2=0$
	\item $\lim\limits_{k \to \infty}\norm{x(k)-x}_{\infty}=0$ 
\end{enumerate}
\minisec{Beweis:}
\underline{$1 \to 2$:} 
\[
	\norm{x(k)-x}_1:=\sum^{n}_{i=1}\underset{\stackrel{\text{(wg 1)}}{\to}0}{\underbrace{\abs{x_i(k)-x_i}}} \longrightarrow 0
\]
\\
\underline{$2 \to 3 \to 4$:} Benutze Ungleichungen $(y \in \mathbb{R}^n)$
\[
	\norm{y}_{\infty} \leq \norm{y}_2 \leq \norm{y}_1 \leq n \cdot \norm{y}_{\infty}
\]
also 
\begin{align*}
	&\lim_{k \to \infty}\norm{x(k)-x}_1=0 \\
	&\Rightarrow \lim_{k \to \infty}\norm{x(k)-x}_2=0 \\
	&\Rightarrow \lim_{k \to \infty}\norm{x(k)-x}_{\infty}=0
\end{align*}
\\
\underline{$4 \to 1$:} \begin{gather*}
	\norm{x(k)-x}_{\infty} = \max\{ \abs{x_i(k)-x} \,| i=1,\dots,n \}
	\Rightarrow \text{für alle $i$: } \abs{x_i(k)-x_i} \leq \norm{x(k)-x}_{\infty}
\end{gather*}
Also
\[
	\left[ \norm{x(k)-x}_{\infty} \stackrel{k \to \infty}{\longrightarrow}0 \right] \qquad \Rightarrow \qquad \left[ \forall\, i=1,\dots,n \quad \abs{x_i(k)-x_i} \longrightarrow 0 \right]
\]
$\Rightarrow 1)$ \bewende

\minisec{Bemerkung:} $X$ metrischer Raum $(x_n)_n , (y_n)_n$ Folgen in $X \quad \lim_{n \to \infty}x_n=x \quad \text{und} \quad \lim_{n \to \infty}y_n=y$
\[
	\Rightarrow d(x,y)=\lim_{n \to \infty}d(x_n,y_n) \qquad \text{(Blatt 2, Aufgabe 2)}
\]
Daraus folgt: \\ Ist $(x_n)_n $ Folge in $X$ und gilt $\lim_{n \to \infty}x_n=x$ sowie $\lim_{n \to \infty}x_n=x'$ dann ist $x=x'$ \\
Denn \[
	d(x,x')=\lim_{n \to \infty}\underset{=0}{\underbrace{d(x_n,x_n)}}=0
\]
\[
	\stackrel{\text{(M2)}}{\Rightarrow} x=x'
\] \bewende
% subsubsection beispiel (end)

\subsubsection[Cauchyfolge im metrischen Raum]{Definition} % (fold)
\label{ssub:definition}
Sei $(X,d)$ metrischer Raum und $(x_n)_n $ Folge in $X$. Dann heißt $(x_n)_n$ {\bfseries Cauchyfolge}, falls gilt:
\[
	\forall\, \varepsilon >0 \, \exists \, n_0 \, \forall\, n,m \geq n_0: \,\, d(x_n,x_m)<\varepsilon 
\]

\minisec{Bemerkung:} $(X,d)$ metrischer Raum $(x_n)_n $ Folgen in $X$, $\lim_{n \to \infty}x_n=x \in X$. \\
Dann ist $(x_n)_n$ Cauchyfolge, denn \\
Sei $\varepsilon >0$. Wähle $n_0$, sodass \[  \forall\, n \geq n_0 \, d(x,x_n)< \frac{\varepsilon}{2} \]
Falls $n,m \geq n_0$
\[
	d(x_n,y_n) \leq \underset{<\frac{\varepsilon}{2}}{\underbrace{d(x_n,x)}} + \underset{<\frac{\varepsilon}{2}}{\underbrace{d(x,y_n)}} < \varepsilon 
\]
% subsubsection definition (end)

\subsubsection[Vollständigkeit]{Definition} % (fold)
\label{ssub:definition}
$(X,d)$ metrischer Raum. Dann heißt $X$ {\bfseries vollständig}, falls jede Cauchy-Folge $(x_n)_n $ in $X$ gegen ein $x \in X$ konvergiert (Ana I)

\minisec{Bsp:} 

	$(0,1) \subseteq \mathbb{R}$.
	$x_n = \frac{1}{n}$ Folge in $X$ ist Cauchyfolge. \\
	$(x_n)_n$ konvergiert \underline{nicht} in $X$, denn der einzig möglicher Grenzwert (0) liegt nicht in $X$.

% subsubsection definition (end)

\subsubsection{Satz} % (fold)
\label{ssub:satz}
$\mathbb{R}^n$, mit der durch die Norm $\norm{.}_p \qquad p=1,2,\infty$ induzierten Metrik ist vollständig

\minisec{Beweis:}
\underline{$p=\infty$:} Sei $x(k)$ $l^{\infty}$-Cauchyfolge in $\mathbb{R}^n$. Dann gilt für $i=1,\dots,n$:
\[
	\abs{x_i(k)-x_i(l)} \leq \norm{x(k)-x(l)}_{\infty}
\]
Da $(x(k))_k$ $l^{\infty}$-Cauchyfolge, ist $(x_i(k))_k$ eine Cauchyfolge in $\mathbb{R}$ für alle $1\leq i \leq n$ \\
Nach dem Cauchy-Kriterium aus Analysis I sind die Folgen $(x_i(k))_k$ Konvergent. $x_i=\lim_{k \to \infty}x_i(k) \in \mathbb{R}$
\[
	x := \sum_{i=1}^{n}x_ie_i \in \mathbb{R}^n
\] 
Dann gilt 
\[
	\norm{x(k)-x}_{\infty} = \max\{ \underset{\to 0}{\underbrace{\abs{x_i(k)-x_i}}} | i=1,\dots,n \} \longrightarrow 0
\]
also
\[
	\lim_{k \to \infty}\norm{x(k)-x}_{\infty} = 0
\]
\\
\underline{$p=2$:}
$x(k)$ $l^2$-Cauchyfolge, weil $\norm{x(k)-x(l)}_{\infty} \leq \norm{x(k)-x(l)}_2$ \\
Ist $(x_k)_k$ eine $l^{\infty}$-Cauchfolge, so gibt es ein $x \in \mathbb{R}^n$ mit 
\[
	\lim_{k \to \infty}\norm{x(k)-x}_{\infty}=0
\]
\[
	\Rightarrow \norm{x(k)-x}_2 \leq n \cdot \underset{\to 0}{\underbrace{\norm{x(k)-x}_{\infty}}} \longrightarrow 0
\] 
\\
\underline{$p=1$:} Analog zum Fall $p=1$, benutze $\norm{y}_{\infty} \leq \norm{y}_1 \leq n \cdot \norm{y}_{\infty}$ 
\bewende	
% subsubsection satz (end)
\subsubsection[Stetigkeit]{Definition} % (fold)
\label{ssub:definition}
Seien $(X,d_X)$ und $(Y,d_Y)$ metrische Räume und $f:X \to Y$ eine Abbildung sowie $x \in X$.
\begin{enumerate}[(i)]
	\item $f$ heißt {\bfseries $\varepsilon$-$\delta$-stetig} (oder {\bfseries stetig}) bei $x$, wenn gilt
	\[
		\forall\, \varepsilon >0 \, \exists \, \delta>0 \, \forall\, x' \in X \text{ mit } d_X(x,x')< \delta \text{ gilt: } d_Y(f(x),f(x'))<\varepsilon 
	\]
	\item $f$ heißt {\bfseries folgenstetig} in $x$ $:\Leftrightarrow$ Ist $(x_n)_n $ eine Folge in $X$ mit $\lim_{n \to \infty}x_n=x$, so gilt
	\[
		\lim_{n \to \infty}f(x_n)=f(x)
	\]
\end{enumerate}
$f$ ist stetig (folgenstetig) auf $X$, wenn $f$ stetig (folgenstetig) bei $x$ ist für alle $x \in X$.
% subsubsection definition (end)

\subsubsection{Satz} % (fold)
\label{ssub:satz}
$f$ stetig in x $\Leftrightarrow$ $f$ folgenstetig in $x$
\minisec{Beweis} wie in Analysis I, wird bald nachgeliefert

% subsubsection satz (end)
\subsubsection{Satz} % (fold)
\label{ssub:satz}
$(X,d_X)$ , $(Y,d_Y)$, $(Z,d_Z)$ metrische Räume. 
\[ f:X \to Y \qquad g:Y \to Z \qquad x \in X \]
Falls $f$ stetig bei $x$ und stetig bei $f(x)$ so ist $g \circ f$ stetig bei $x$.
\minisec{Beweis:} Siehe Analysis I 
% subsubsection satz (end)
% subsection euklidsche_und_normierte_vektorraume (end)
\subsubsection{Beispiele} % (fold)
\label{ssub:beispiele}
\begin{align*}
&\mu: \mathbb{R}^2 \to \mathbb{R} \qquad \mu(x,y):= \, xy \qquad \text{(Multiplikation)} \\
&\alpha: \mathbb{R}^2 \to \mathbb{R} \qquad \alpha(x,y):=x+y \qquad \text{(Addition)}  \\
&q:\mathbb{R} \times \mathbb{R} \setminus {0} \to \mathbb{R} \qquad q(x,y):= \frac{x}{y} \qquad \text{(Division)}
\end{align*}
sind auf ihrem gesamten Definitionsbereich stetig.

\minisec{Beweis:}
\underline{z.\,B $q$:}
$q$ ist folgenstetig, denn: \\
Sei $(x,y) \in \mathbb{R} \times \mathbb{R} \setminus \{0\}$ (d.h $y \neq 0$) und $(x_n,y_n)$ Folge in $\mathbb{R} \times \mathbb{R} \setminus \{0\}$ mit 
\[
	\lim_{n \to \infty}(x_n,y_n)=(x,y)
\]
Dann gilt: 
\[
	\lim_{n \to \infty}x_n=x \text{ und } \lim_{n \to \infty}y_n=y \neq 0
\]
\[
	\Rightarrow q(x_n,y_n)= \frac{x_n}{y_n} \stackrel{\text{Grenzwertsätze}}{\longrightarrow } \frac{x}{y} = q(x,y)
\]
\[
	\Rightarrow q \text{ ist folgenstetig}
\]
$\alpha, \mu$ analog.
\bewende

 \minisec{\underline{Bsp:}}
 Sei $X$ metrischer Raum, $f,g:X \to \mathbb{R}$ stetig \\
 $\Rightarrow $ 
 Die Funktion 
 \[
 \begin{array}{rccl}
 	(f,g): &X &\to &\mathbb{R}^2 \\
	&x &\mapsto &(f(x),g(x))\in \mathbb{R}^2
 \end{array}
 \]
 ist stetig. \\
 Denn, $(x_n)_n $ Folge in $X$, $x_n \to x \in X$ \\ 
 $$\Rightarrow f(x_n) \to f(x) \text{ und } g(x_n) \to g(x) \qquad (f,g \text{ stetig})$$
 \[
 	\stackrel{\text{Tatsache}}{\Rightarrow } (f,g)(x_n)=(f(x_n),g(x_n)) \to (f(x),g(x)) \in \mathbb{R}^2 \\
 \]
 \[
	\Rightarrow (f,g) \text{ stetig}
 \]
 Betrachte Kompositionen 
 \[
 \begin{array}{rcccl}
 	X &\stackrel{(f,g)}{\to }  &\mathbb{R}^2 &\stackrel{\mu}{\to } &\mathbb{R} \\
	x &\stackrel{\hphantom{(f,g)}}{\mapsto } &(f(x),g(x)) &\mapsto &f(x)g(x) \\
 \end{array}
 \]
 \[
 	\Rightarrow \underset{\text{$\Rightarrow$ stetig}}{\underbrace{ \underset{\text{stetig}}{\underbrace{\mu}} \circ \underset{\text{stetig}}{\underbrace{(f,g)}} }}
	= fg
 \] 
 \[
 	\Rightarrow \text{Produkte stetiger Funktionen sind wieder stetig}
 \]
 \minisec{Analog} Summen Quotienten stetiger Funktionen sind stetig
% subsubsection beispiele (end)
\subsubsection[Lipschitz-stetig]{Definition} % (fold)
\label{ssub:definition}
$(X,d_X),(Y,d_y)$ metrische Räume $f:X \to Y$ heißt {\bfseries Lipschitz-stetig} 
$:\Leftrightarrow \exists \, L \leq 0, \text{ so dass gilt: }$
\begin{align*}
	 \forall\, x,x' \in X : \qquad 
	 d_Y(f(x),f(x')) \leq Ld_X(x,x')
\end{align*}
\minisec{Bem:} ist $f$ Lipschitz-stetig, so ist $f$ stetig: \\
sei $\varepsilon > 0$ Wähle $\delta$, so dass $\delta\,L \leq \varepsilon$ \\
Dann:
\[
	\text{Ist } d_X(x,x')<\delta \Rightarrow d_Y(f(x),f(x')) \leq L \cdot d_X(x,x') < \delta L \leq \varepsilon \qquad \Rightarrow f \text{ stetig.} 
\]
% subsubsection definition (end)

\subsubsection[Beispiel Lipschitz-stetig]{Beispiel} % (fold)
\label{ssub:beispiel}
Betrachte $\mathbb{R}^n$ mit der durch die $l^p$-Normen induzierten Metrik $p \in \{1,2,\infty\}$ \\
Dann gilt: Sind $p,q \in \{1,2,\infty\}$, so ist 
\[
	\text{id:}(\mathbb{R}^n,\norm{.}_p) \to (\mathbb{R}^n,\norm{.}_q) 
\] 
Lipschitz-stetig.
\minisec{Grund:} $\forall\, x \in \mathbb{R}^n$
\[
	\norm{x}_{\infty} \leq \norm{x}_2 \leq \norm{x}_1 \leq n \cdot \norm{x}_{\infty}
\]
% subsubsection beispiel (end)
\subsubsection{Satz} % (fold)
\label{ssub:satz}
Jede lineare Abbildung $f:\mathbb{R}^n \to \mathbb{R}^m$ ist Lipschitz-stetig.

\minisec{Beweis:} Aus LA I ist bekannt: \\
Es gibt genau ein $A \in \ \text{Mat}_{m,n}(\mathbb{R})$, so dass $f(x)=Ax \qquad (\forall\, x \in \mathbb{R}^n)$ 
\[
	\text{M}_{m,n}(\mathbb{R}) = \text{Mat}_{m,n}(\mathbb{R})
\]
% subsubsection satz (end)
\minisec{Beh:}
$\exists \, L=L_A$, so dass für alle $x \in \mathbb{R}^n$ gilt 
\[
	\norm{Ax}_2 \leq L \norm{x}_2
\]
Es gilt 
\[
	\norm{Ax}_2 \leq \norm{Ax}_1 = \sum^{m}_{i=1} \abs{(Ax)_i} = \sum^{m}_{i=1} \left|\sum^{n}_{j=1} a_{ij}x_j \right| 
\]
nach Cauchy-Schwarz gilt:
\[
	\left| \sum^{n}_{j=1}y_jz_j \right | =\abs{\skal{y}{z}} \stackrel{\text{CS}}{\leq } \norm{y}_2 \cdot \norm{z}_2 = \left( \sum^{m}_{j=1}y_j^2 \right)^{\frac{1}{2}}
	 \left( \sum^{m}_{j=1}z_j^2 \right)^{\frac{1}{2}}
\]
Also 
\[
	\Rightarrow \left| \sum^{m}_{j=1}a_{ij}x_j \right| \leq \left( \sum^{n}_{j=1}a_{ij}^2 \right)^{\frac{1}{2}} \underset{\norm{x}_2}{\underbrace{
	\left ( \sum^{n}_{j=1}x_{j}^2 \right ) }}^{\frac{1}{2}}
\]
\[
	\Rightarrow \norm{Ax}_2 \leq \sum^{n}_{i=1} \left( \left( \sum^{n}_{j=1}a_{ij}^2y \right)^{\frac{1}{2}}\norm{x}^2 \right) = \underset{=L}{\underbrace{\sum^{n}_{i=1} \left( \sum^{n}_{j=1}a_{ij}^2 \right)^{\frac{1}{2}}}} \norm{x}_2
\] 
\bewende
\subsubsection{Korollar:} % (fold)
\label{ssub:korollar}
$A \in \text{Mat}_{m,n}(\mathbb{R})$ \\
Dann ist 
\[
	\norm{A}:= \sup\limits_{\norm{x}_2 \leq 1} \norm{Ax}_2 < \infty
\]
\minisec{Beweis:}
Ist $L \geq 0$ mit $\norm{Ax}_2 \leq L \cdot \norm{x}_2 \qquad \forall\, x \in \mathbb{R}^n$ \\ 
so gilt für $\norm{x}\leq 1$: \, $\norm{Ax}_2 \leq L$
\[
	\Rightarrow \sup\limits_{\norm{x}\leq 1}\norm{Ax}_2 \leq L < \infty
\]
 \bewende
% subsubsection korollar (end)
\subsubsection[Operatornorm]{Definition} % (fold)
\label{ssub:definition}
Die Zahl $\norm{A}$ heißt {\bfseries Operatornorm} von $A$ \\
Es gilt: Die Operatornorm ist eine Norm auf dem $\mathbb{R}$-VR \, $\text{Mat}_{m,n}(\mathbb{R})$ \qquad (Übung) \\
Es gibt keine einfache Formel für $\norm{A}$
% subsubsection definition (end)

\subsection{Topologische Grundbegriffe} % (fold)
\label{sub:topologische_grundbegriffe}
\subsubsection[$\varepsilon$-Ball]{Definition} % (fold)
\label{ssub:definition}
 $(X,d)$ sei metrischer Raum, $\varepsilon >0$ und $x \in X$. Der {\bfseries $\varepsilon $-Ball} in $X$ um $x$ ist
 \[
 	B_{\varepsilon }(x) = \{ y \in X | d(x,y)< \varepsilon \}
 \]
 in $\mathbb{R}$: \[
 	B_{\varepsilon}(x)=(x-\varepsilon ,x+ \varepsilon)
 \]
% subsubsection definition (end)
\subsubsection[Innerer Punkt, Umgebung]{Definition} % (fold)
\label{ssub:definition}
$(X,D)$ metrischer Raum $Y \subseteq X, x \in X$
\begin{enumerate}
	\item $x$ heißt {\bfseries innerer Punkt} von $Y$  \[
		\exists\, \delta>0 \text{ mit } B_{\delta }(x) \subseteq Y
	\]
	In diesem Fall heißt eine {\bfseries Umgebung} von $x$ in $X$
	\item $x \in X$ heißt {\bfseries Berührpunkt} von $Y$
	\[
		\forall\, \varepsilon >0 \, \exists\, y \in Y \text{ mit } d(x,y)<\varepsilon 
	\]
	(Äquivalent: $\forall\, \varepsilon > 0 \qquad B_{\varepsilon }(x)\cap \neq \varnothing$ )
	\begin{align*}
		\bar{y} &= \{ x \in X \,|\, x \text{ innerer Punkt von } Y \} \\
		\overset{\circ}{y} &= \{ x \in X \,|\, x \text{ ist Berührpunkt von Y} Y \} 
	\end{align*}
\end{enumerate}
% subsubsection definition (end)
\subsubsection[Offen, Abgeschlossen]{Definition} % (fold)
\label{ssub:definition}
\begin{enumerate}
	\item $(X,d)$ metrischer Raum, $Y \subseteq X$. Dann heißt {\bfseries $Y$ offen} in $X$ $\Leftrightarrow $ jedes $x \in Y$ ist innerer Punkt von $Y$ \\
	( $\Leftrightarrow \forall\, x \in Y \, \exists \, \delta >0: \, B_{\delta } \subseteq Y$)
	\item $Y$ heißt {\bfseries abgeschlossen} in $X$ $:\Leftrightarrow $ Das Komplement $X-Y$ ist offen in $X$
\end{enumerate}
 \minisec{Bsp:}
 $X = \mathbb{R}$
 \begin{enumerate}
 	\item Offene Intervalle $(a,b) \subseteq \mathbb{R}$ sind offen in $\mathbb{R}$,denn
  \begin{align*}
  	&x \in (a,b) \,,\, \delta = \min \{ x-a.b-x \} >0 \\
	&\Rightarrow (x-\delta ,x+\delta ) \subseteq (a,b)
  \end{align*}  
  
 \item Abgeschlossene Intervalle $[a,b] \subseteq \mathbb{R}$ sind abgeschlossen: 
 \begin{align*}
 	&\mathbb{R} \setminus [a,b] = (-\infty,a) \cup (b, +\infty) \\
 	&\Rightarrow (x-\delta ,x+\delta ) \subseteq \mathbb{R} \setminus [a,b] \\
	&\Rightarrow x>b \text{ innerer Punkt von } \mathbb{R} \setminus[a,b] \\
	&\quad \text{ analog } x<a \text{ ist innerer Punkt von } \mathbb{R}\setminus[a,b] \\
	&\Rightarrow \mathbb{R}\setminus[a,b] \text{ offen in } \mathbb{R} \Rightarrow  [a,b] \text{ abgeschlossen in } \mathbb{R} 
\end{align*}
\end{enumerate}
% subsubsection definition (end)

\minisec{Warnung:}
Offenheit ist eine Eigenschaft von \underline{Teilmengen} metrischer Räum, nicht der Mengen selber. '$Y$ offen' sinnlos ! '$Y$ offen in $X$' sinnvoll! Analog für Abgeschlossen
 \\
$X = \mathbb{R}^2$ \\
\begin{align*}
	&Z = \{ (x,y) \in \mathbb{R}^2\, |\, y=0 \} \\
	&Y = \{ (x,y) \in \mathbb{R}^2 \, | \, y=0,0<x<1 \} \\
	&Y \text{ offen in } Z \\
	&Y \text{ nicht offen in } X \\
	\end{align*}
% subsection topologische_grundbegriffe (end)
\subsubsection{Beispiel} % (fold)
\label{ssub:beispiel}
$X$ metrischer Raum, $ x \in X$,$\varepsilon>0$. Dann ist $B_{\varepsilon}^X(x)$ offen in $X$

\minisec{Beweis:}
zu zeigen: für jedes $y \in B_{\varepsilon}(x) \exists\, \delta >0: \, B_{\delta }(y) \subseteq B_{\varepsilon}(x)$
\[
	d(x,y)<\varepsilon \Rightarrow  \exists\, \delta>0: \, d(x,y)+\delta <\varepsilon
\] \\
Beh: $B_{\delta }(y) \subseteq B_{\varepsilon}(x)$ \\
Das ist wahr, denn
\[
	z \in B_{\delta}(y) \Rightarrow d(y,z)<\delta 
\] \\
Ferner: 
\[
	d(z,x) \leq d(z,y)+d(y,x)<\delta + d(y,x)<\varepsilon, \text{d.h.} z \in B_{\varepsilon}(x)
\]
% subsubsection beispiel (end)
\subsubsection{Satz} % (fold)
\label{ssub:satz}
Es sei $(X,d)$ metrischer Raum. \\
Dann gilt:
\begin{enumerate}
	\item $\varnothing$ und $X$ sind offen in $X$
	\item Falls $U_0,U_1$ offen in $X$ sind, so auch $U_0 \cap U_1$
	\item Falls $(U_i)_{i \in I}$ eine durch $I$ induzierte Familie von Teilmengen von $X$ ist, so ist $\bigcup\limits_{i \in I}U_i$ offen in $X$ \\
		'beliebige Vereinigungen offener Mengen sind offen'
	\item $\varnothing,X$ sind abgeschlossen in $X$
	\item $A_0,A_1$ abgeschlossen in $X$, dann ist $A_0 \cup A_1$ abgeschlossen in $X$
	\item $(A_i)_{i \in I}$, $A_i$ abgeschlossen in $X$, so ist $\bigcap\limits_{i \in I}A_i$ abgeschlossen in $X$
\end{enumerate}
\minisec{Beweis:}
\begin{enumerate}
	\item $\varnothing \subseteq X$ offen trivial \\
		$X \subseteq  X$ offen trivial
	\item Sei $x \in U_0 \cap U_1$ Dann existieren Bälle
	\begin{align*}
		B_{\varepsilon_1}^X(x) \subseteq U_1 \qquad B_{\varepsilon_0}^X(x) \subseteq U_0
	\end{align*}
Wähle nun $\varepsilon= \min \{ \varepsilon_0, \varepsilon_1 \} > 0$
\begin{align*}
	& \Rightarrow B_{\varepsilon}(x) \subseteq U_0 \cap U_1 \\
	& \Rightarrow x \text{ innerer Punkt von } U_0 \cap U_1
\end{align*}
	\item Sei $x \in \bigcup\limits_{i \in I}U_i \subseteq X$
	\[
		\Rightarrow \exists\,i \in I:\, x \in U_i, \text{ weil } U_i \subseteq X \text{ offen, gibt es } \varepsilon>0: \, B_{\varepsilon}(x) \in U_i.
	\]
Aber \[
	U_i \subseteq \bigcup\limits_{i \in I}U_i, \text{ daher: } B_{\varepsilon}(x) \subseteq \bigcup\limits_{i \in I}U_i
	\]
	\item \begin{align*}
		&\varnothing = X \setminus \underset{\text{offen nach 1)}}{\underbrace{X}} \qquad \Rightarrow \varnothing \text{ abgeschlossen in } X \\
		&X = X \setminus \underset{\text{offen}}{\underbrace{\varnothing}} \qquad \Rightarrow X \subseteq X \text{ abgeschlossen}
	\end{align*}
	\item \[
		X \setminus(A_0 \cup A_1)= \underset{\text{offen in $X$}}{\underbrace{\underset{\text{offen in $X$}}{\underbrace{(X \setminus A_0)}}\cap \underset{\text{offen in $X$}}{\underbrace{(X \setminus A_1)}}}}
	\]
	\[
		\Rightarrow A_0 \cup A_1 \subseteq X \text{ abgeschlossen}
	\]
	\item \[
		X \setminus \bigcap\limits_{i \in I}A_i = \underset{\text{offen nach 3)}}{\underbrace{\bigcup\limits_{i \in I}\underset{\text{offen in $X$}}{\underbrace{X \setminus A_i}}}}
	\]
	\[
		\Rightarrow \bigcap\limits_{i \in I}A_i \text{ abgeschlossen in } X
	\]
\end{enumerate}
% subsubsection satz (end)
\minisec{Bemerkung:}
Beliebige Durchschnitte offener Mengen sind i.A. nicht offen 
\minisec{\underline{Beispiel}:}
\[
	U_n = (-\frac{1}{n},\frac{1}{n}) \text{ ist offen in } \mathbb{R}
\]
\[
	\bigcap\limits_{n \in \mathbb{N}}U_n = \{0\} \text{ ist nicht offen in } \mathbb{R}
\]

\subsubsection{Satz} % (fold)
\label{sub:Satz}
$X$ metrischer Raum, $Y \subseteq X$. Dann sind Äquivalent:
\begin{enumerate}
	\item $Y \subseteq X$ abgeschlossen in $X$
	\item Für jede Folge $(y_n)$ mit $y_n \subseteq Y$, welche in $X$ konvergiert, gilt
	\[
		\lim_{n \to \infty}y_n \in Y
	\]
	\end{enumerate}
\minisec{Beweis:}
\minisec{Vorbemerkung:} $y \in X$ ist genau dann ein Berührpunkt von $Y \subseteq X$, wenn es eine Folge $(y_n)_n$ gibt mit $(y_n) \in Y$ und $\lim_{n \to \infty}y_n = y$
\minisec{Beweis der Vorbemerkung:}
$"\Leftarrow"$: Ist $\lim_{n \to \infty}y_n=y \qquad y_n \in Y$
\begin{align*}
	& \Rightarrow \forall\, \varepsilon>0 \,\exists\,n: \, d(y,y_n)<\varepsilon \\
	&\Rightarrow \forall\, \varepsilon>0 \, y_n \in B_{\varepsilon}(y) \cap Y \neq \varnothing \\
	&\Rightarrow y \text{ ist Berührpunkt von } Y
\end{align*}
$"\Rightarrow"$:
\begin{align*}
	&y \text{ Berührpunkt von } Y \\
	&\Rightarrow  \forall\, \varepsilon > 0 \, \exists\, y' \in Y: \, d(y,y')<\varepsilon
\end{align*}
Wähle für $n \in \mathbb{N}$ ein $y_n \in Y$ mit $d(y,y_n)<\frac{1}{n}$
\[
	\Rightarrow \lim_{n \to \infty}\underset{\in y}{\underbrace{y_n}}=y
\] \bewende
\minisec{Beweis des Satzes:}
Sei 1) nicht erfüllt. Dann:
\begin{align*}
	& \stackrel{\hphantom{\text{Vorbem.}}}{\Leftrightarrow}  Y \subseteq X \text{ nicht abgeschlossen in } X \\
	& \stackrel{\hphantom{\text{Vorbem.}}}{\Leftrightarrow} X \setminus Y \text{ nicht offen in } X \\
	& \stackrel{\hphantom{\text{Vorbem.}}}{\Leftrightarrow} \exists\,y \in X \setminus Y, \, y \text{ kein innerer Punkt von } X \setminus Y \\
	& \stackrel{\hphantom{\text{Vorbem.}}}{\Leftrightarrow} \exists\,y \in X \setminus Y, \, \forall\, \varepsilon>0 \qquad B_{\varepsilon}(y) \not\subseteq X
	 \setminus Y
	 \\
	& \stackrel{\hphantom{\text{Vorbem.}}}{\Leftrightarrow} \exists\,y \in X \setminus Y, \, \forall\, \varepsilon>0 \qquad B_{\varepsilon}(y) \cap Y \neq \varnothing
	 \\
	& \stackrel{\hphantom{\text{Vorbem.}}}{\Leftrightarrow} \exists\,y \in X \setminus Y, \, y \text{ ist Berührpunkt von } Y \\
	& \stackrel{\text{Vorbem.}}{\Leftrightarrow} \exists\, \text{ Folge } y_n \text{ in } Y, y_n \to y \notin  Y
\end{align*}
und dazu ist Äquvivalent, dass 2) nicht erfüllt ist.
\bewende
% subsection Satz (end)
\subsubsection{Satz} % (fold)
\label{ssub:satz}
$(X,d_X)$,$(Y,d_Y)$ metrischer Räume, $f: X \to Y$ Abbildung. Dann sind Äquivalent:
\begin{enumerate}
	\item $f$ ist folgenstetig.
	\item Für jede abgeschlossene Teilmenge $A \subseteq Y$ ist $f^{-1}(A) \subseteq X$ abgeschlossen.
	\item Für jede offene Teilmenge $U \in Y$ ist $f^{-1}(U)$ offen in $X$
	\item $f$ ist $\varepsilon$-$\delta$-stetig
\end{enumerate}
\minisec{Beweis:}
$\underline{1 \Rightarrow 2}$: Sei $f$ folgenstetig, $A \subseteq Y$ abgeschlossen. \\
zu zeigen: (wegen I.4.7) ist $x_n \in f^{-1}(A)$ und $x_n \to x$, so gilt $x \in f^{-1}(A)$ \\
\begin{align*}
	&x_n \to x \Rightarrow \underset{\in A}{\underbrace{f(x_n)}} \to \underset{\in y}{\underbrace{f(x)}} \qquad \text{($f$ ist folgenstetig)} \\
	&A \subseteq Y \text{ abgeschlossen } \Rightarrow  f(x) \in A \\
 	&\Rightarrow x \in f^{-1}(A) 
\end{align*}
$\underline{2 \Rightarrow 3}$:
\begin{align*}
	&U \subseteq Y \text{ offen } \Rightarrow Y \setminus U \subseteq Y \text{ abgeschlossen} \\
	&\stackrel{2)}{\Rightarrow} f^{-1}(Y \setminus U) \subseteq  X \text{ abgeschlossen.} \\
	&f^{-1}(Y \setminus U) = X \setminus f^{-1}(U) \Rightarrow f^{-1}(U) \subseteq  X \text{ offen}
\end{align*}
$\underline{3 \Rightarrow 4}$:
Sei $x \in X$, $\varepsilon>0$ 
\begin{align*}
	&\Rightarrow B_{\varepsilon}^Y(f(x)) \subseteq Y \text{ offen in } Y \\
	&\stackrel{3)}{\Rightarrow} f^{-1}(B_{\varepsilon}^Y(f(x))) \subseteq X \text{ offen} \\
	&\stackrel{\text{def. offen}}{\Rightarrow} \exists\, \delta >0 \, B_{\delta}^X(x) \subseteq f^{-1}(B_{\varepsilon}^Y(f(x)))
\end{align*}
Das heißt: ist $x' \in X$ mit $d(x,x')<\delta$
\[
	\Rightarrow d(f(x),f(x'))<\varepsilon
\]
$\underline{4 \Rightarrow 1}$:
Sei $f$ $\varepsilon$-$\delta$-stetig und $(x_n)_n$ Folge in $x$. $x_n \to x \in X$. Dann gilt $f(x_n) \to f(x)$, denn: \\
Sei $\varepsilon>0$
\begin{align*}
	&\Rightarrow \exists\,\delta >0: x' \in X \, d(x,x')< \delta \\
	&\Rightarrow d(f(x),f(x'))<\varepsilon \Rightarrow \text{weil } x_n \to x: \, \exists\,n_0: \, \forall\, n \geq n_0: \, d(x,x_n)<\delta \\
	&\Rightarrow \exists\,n_0: \, \forall\, n \geq n_0: \, d(f(x),f(x_n))<\varepsilon \text{ d.h. } f(x_n) \to f(x)
\end{align*}
% subsubsection satz (end)
\subsubsection{Beispiel} % (fold)
\label{ssub:beispiel}
$(V,\norm{.})$ normierter Vektorraum
\[
	\norm{.}: V \to \mathbb{R} \text{ stetig, denn } \norm{x}-\norm{y} \leq \norm{x-y}
\]
Also sind die Mengen (für $x \in V$, $r \geq 0$)
\[
	\{y \in V \, | \, \norm{x-y}\leq r \} \subseteq V \text{ und } \{y \in V \,|\, \norm{x-y}=r \} \subseteq V
\] abgeschlossen.
% subsubsection beispiel (end)
\subsection{Kompaktheit} % (fold)
\label{sub:kompaktheit}

\minisec{Erinnerung: Satz von Bolzano-Weierstraß}
'Jede beschränkte Folge in $\mathbb{R}$ besitzt eine konvergente Teilfolge'
\begin{itemize}
	\item $\rightarrow $ Cauchy-Konvergenzkriterium
	\item $\rightarrow $ Satz vom Minimum und Maximum (Wenn $f:A \to B$ stetig, dann gilt: $f$ nimmt Minimum und Maximum an)
	\item $\to$ Mittelwertsatz der Differentialrechnung
	\item $\to$ \underline{Satz:} $f:[a,b] \to \mathbb{R}$ stetig $\Rightarrow$ f ist gleichmäßig stetig. \\
		  $\Rightarrow$ stetige Funktionen sind Riemann-integrierbar
\end{itemize}
\subsubsection{Satz von Bolzano-Weierstraß im $\mathbb{R}^n$} % (fold)
\label{ssub:satz_von_bolzano_weierstrass_im_mathbb_r_n}
Jede beschränkte Folge in $\mathbb{R}^n$ besitzt eine konvergente Teilfolge.
 
\minisec{Beweis:}
Sei $(x(k))_k$ Folge in $\mathbb{R}^n$, $\norm{x(k)}_2 \leq C \quad  \forall\, k$. \\
$x_i(k) \quad  i=1,\dots,n$ Komponenten
\[
	\Rightarrow \underset{\text{Folge in } \mathbb{R}}{\underbrace{\abs{x_i(k)}}} \leq \norm{x(k)}^2 \leq C
\]
Nach Bolzano-Weierstraß gibt es eine Teilfolge $(x(k_l))_l$ von $x(k)$, so dass $x_1(k_l)$ (in $\mathbb{R}$) Konvergiert
Nenne \[
	x(1,l)=x(k_l)
\]
\[
\begin{matrix}
	\text{1. Komponente} & \text{2. Komponente} & \text{$n$-te Komponente} \\
	x(k_1=x(1,1) & x(2,1) & x(n,1) \\
	x(k_2)=x(1,2) & x(2,2) & x(n,2) \\
	\vdots & \vdots & \vdots 
\end{matrix}
\]
\[
	\begin{matrix}
		\text{Teilfolge} & (x(1,k))_k & (x_1(1,k)_k) \text{ Konvergiert} \\
		\text{Teilfolge} & (x(2,k))_k & (x_2(2,k)_k) \text{ Konvergiert} \\
		\vdots & \vdots & \vdots \\
		\text{Teilfolge} & (x(n,k))_k & (x_n(n,k)_k) \text{ Konvergiert}
	\end{matrix}
\]
\begin{align*}
	&\Rightarrow \text{Alle Komponentenfolgen von } (x(n,k))_k \text{ konvergieren} \\
	&\Rightarrow (x(n,k))_k \text{ ist Konvergent}
\end{align*}
\bewende
% subsubsection satz_von_bolzano_weierstrass_im_mathbb_r_n (end)
\subsubsection{Definition} % (fold)
\label{ssub:definition}
Ein metrischer Raum $X$ heißt {\bfseries folgenkompakt}, falls für jede Folge $(x_n)_n$ in $X$ eine konvergente Teilfolge hat. \\
Eine Teilmenge $Y \subseteq Y$ eines metrischen Raumes heißt folgenkompakt, falls jede Folge in $(y_n)_n$ in $Y$ eine (in $Y$!) konvergente Teilfolge besitzt.
\[
	( \Leftrightarrow Y \text{ mit der von $X$ eingeschränkten Metrik, ist folgenkompakt} )
\]
% subsubsection definition (end)
\subsubsection{Beispiel:} % (fold)
\label{ssub:beispiel}
\begin{enumerate}
	\item $[a,b]$ folgenkompakt
	\item $X= \{ x \in \mathbb{R}^n \,|\, \norm{x}_2 \leq C \}$ ist folgenkompakt: $(x_k)_k$ Folge in $X$
	\[
		\stackrel{I.5.1}{\Rightarrow } \exists\, \text{ TF } x_k \to x \in \mathbb{R}^n
	\]
	aber $X \subseteq \mathbb{R}^n$ ist abgeschlossen, also $x \in X$
	\item $(0,1)$,$(0,1]$ nicht folgenkompakt (betrachte Folge ($\frac{1}{n})_n$)
	\item $\mathbb{R}$ nicht folgenkompakt (Die Folge $(n)_n$ hat keine Konvergente TF)
	\item $\mathbb{Q} \cap [0,1] \subseteq \mathbb{R}$ nicht folgenkompakt.
	\[
		\text{Ist } x \in [0,1], \text{ so gibt es eine Folge } x_n \in [0,1] \cap \mathbb{Q} \text{ mit } x_n \to x
	\]
	Dann hat $x_n$ keine in $[0,1] \cap \mathbb{Q}$ konvergente Teilfolge (falls $x$ irrational)
\end{enumerate}
% subsubsection beispiel (end)

\subsubsection{Satz} % (fold)
\label{ssub:satz}
$X$ metrischer Raum
\begin{enumerate}
	\item Ist $X$ folgenkompakt, so ist $X$ vollständig
	\item $X$ folgenkompakt, $Y \subseteq X$ abgeschlossen $\Rightarrow$ $Y$ folgenkompakt
	\item $Y \subseteq X$, $Y$ folgenkompakt $\Rightarrow$ $Y \subseteq X$ abgeschlossen
\end{enumerate}

\minisec{Beweis:}
\begin{enumerate}
	\item Sei $(x_n)_n$ Cauchyfolge in $X$
	\[
		X \text{ folgenkompakt } \Rightarrow \text{ es gibt konvergente Teilfolge } x_{n_k} \to x \in X 
	\]
	\minisec{Beh.} $x_n \to x$ \\
	Denn \[
		d(x,x_n) \leq d(x,x_{n_k})+d(x_{n_k},x_n)
	\] Sei $\varepsilon>0$ Wähle $n_0 \quad \forall\, n,m \geq n_0: \, d(x_n,x_m)< \frac{\varepsilon}{2}$ \\
	Wähle $k: \quad n_k \geq n_0 $
	\[
		\text{für } n \geq n_k \quad d(x,x_n) \leq \underset{< \frac{\varepsilon}{2}}{\underbrace{d(x,x_{n_k})}} + 
		\underset{< \frac{\varepsilon}{2}}{\underbrace{d(x_{n_k},x)}} y \varepsilon
	\]
	\item Sei $(y_n)_n$ Folge in $Y \subseteq X$
	\[
		X \text{ folgenkompakt } \Rightarrow \exists\, \text{ TF } y_{n_k} \text{, welche gegen ein $x \in X$ konvergiert}
	\]
	Weil $Y \subseteq X$ abgeschlossen, gilt
	\[
		x = \lim_{k \to \infty}y_{n_k} \in Y \qquad \Rightarrow \qquad (y_{n_k})_k \text{ ist in $Y$ konvergent}
	\]
	$\Rightarrow Y$ folgenkompakt
	\item \underline{zu zeigen:} Ist $(y_n)_k$ Folge in $Y$. $y_n \to x \in X \Rightarrow x \in Y$ \\
	$(y_n)$ ist Cauchyfolge. \\
	 Nach 1) ist $Y$ vollständig, dh $\exists\, z \in y: \quad y_n \in z $
	\[
		\Rightarrow x=z \text{, daher } x \in Y
	\]
\end{enumerate}
% subsubsection satz (end)
\subsubsection{Satz von Heine-Borel} % (fold)
\label{ssub:satz}
Eine Teilmenge $X \subseteq \mathbb{R}^n$ ist folgenkompakt $\Leftrightarrow$ $X$ ist beschränkt und abgeschlossen in $\mathbb{R}^n$

\minisec{Beweis:}
\underline{'$\Rightarrow$'}: Sei $X$ folgenkompakt. Dann folgt aus I.5.4 3): $X \subseteq \mathbb{R}^n$ abgeschlossen \\
Wäre $X$ nicht beschränkt, dann gäbe es zu jedem $n \in \mathbb{N}$ ein $x_n \in X$ mit $\norm{x}_2 \geq n$ \\
Die Folge $(x_n)$ kann dann keine konvergente Teilfolge enthalten, also kann $X$ nicht folgenkompakt sein! \\
\\
\underline{'$\Leftarrow$'}: Sei $X \subseteq \mathbb{R}^n$ beschränkt und abgeschlossen \\
Sei $(x_n)_n$  Folge in $X$. \\
Die Folge $(x_n)_n$ ist beschränkt. 
\[
	\stackrel{I.5.1}{\Rightarrow} \text{ es gibt konvergente Teilfolge } x_{n_k} \to x \in \mathbb{R}^n
\]
Weil $X$ abgeschlossen in $\mathbb{R}^n$, ist $x \in X$
\[
	\Rightarrow X \text{ folgenkompakt}
\]
\bewende

\minisec{Beispiele:}
\begin{itemize}
	\item $D^n = \{ x \in \mathbb{R}^n \, | \, \norm{x} \leq 1 \} \subseteq \mathbb{R}^n$ \qquad \underline{abgeschlossene Einheitskugel}
	\item $S^{n-1}= \{ x \in \mathbb{R}^n \, | \, \norm{x}_2=1 \} \subseteq \mathbb{R}^n$ \qquad \underline{Einheitssphäre}
\end{itemize}
$\Rightarrow$ $D^n$,$S^{n-1}$ folgenkompakt, denn $D^n$,$S^{n-1}$ beschränkt und:
\[
	f: \mathbb{R}^n \to \mathbb{R} \qquad f(x)=\norm{x}_2 \qquad f \text{ stetig}
\]
\[
	D^n = f^{-1}(\underset{\text{abg. in $\mathbb{R}$}}{\underbrace{[0,1]}}) \qquad S^{n-1}=f^{-1}(\underset{\text{abg. in $\mathbb{R}$}}{\underbrace{\{1\}}})
\]
$\Rightarrow$ $D^n$,$S^{n-1}$ abgeschlossen
% subsubsection satz (end)

\subsubsection{Satz} % (fold)
\label{ssub:satz}
$X$,$Y$ metrische Räume, $f: X \to Y$ stetig. Sei $X$ folgenkompakt 
\[
	\Rightarrow f(X) \subseteq Y \text{ ist folgenkompakt}
\]
\minisec{Beweis:}
Sei $y_n$ Folge in $f(X)$ \\
\[ \forall\, n \, \exists\, x_n \in X \text{ mit } f(x_n)=y_n \] 
$(x_n)_n$ Folge in $X$
\[
	X \text{ folgenkompakt } \Rightarrow  \exists\, \text{ konvergente Teilfolge } x_{n_k} \to x \in X
\]
\[
	f \text{ stetig } \Rightarrow y_{n_k} = f(x_{n_k}) \to f(x) \in f(X)
\]
Damit folgt die Behauptung
\bewende

% subsubsection satz (end)
\subsubsection{Satz vom Minimum und Maximum} % (fold)
\label{ssub:satz_vom_minimum_und_maximum}
$X$ folgenkompakt, $f: X \to \mathbb{R}$ stetig. \\ Dann existieren $x_{\text{min}},x_{\text{max}} \in X$, so dass für alle $x \in X$ gilt:
\[
	f(x_{\text{min}}) \leq f(x) \leq x_{\text{max}}
\]
\minisec{Beweis:}
$f(X) \subseteq \mathbb{R}$ ist folgenkompakt (Satz I.5.6) \\
Nach I.5.5 ist $f(X) \subseteq \mathbb{R}$ abgeschlossen und beschränkt. \\
Seien $c,C \in \mathbb{R}$ definiert durch
\begin{align*}
	c &:= \inf f(x) \\
	C &:= \sup f(x)
\end{align*}
Dann sind nach Ana I $c,C$ Berührpunkte um $f(x)$
\[
	f(x) \text{ abgeschlossen } \Rightarrow c,C \in f(x)
\]
Wähle nun $x_{\text{min}},x_{\text{max}} \in X$ mit
\[
	f(x_{\text{min}})=c \,,\,f(x_{\text{max}})=C
\] \bewende
% subsubsection satz_vom_minimum_und_maximum (end)
% subsection kompaktheit (end)
\end{document}	